\chapter{2024年天津卷}

\section{单选题}

\begin{problem}
集合 \(A=\{1,2,3,4\}\)，\(B=\{2,3,4,5\}\)，则 \(A\cap B=\)（\quad）

\choices
    {\(\{1,2,3,4\}\)}
    {\(\{2,3,4\}\)}
    {\(\{2,4\}\)}
    {\(\{1\}\)}
\end{problem}

\begin{problem}
设 \(a\)，\(b\in \R\)，则“\(a^3=b^3\)”是“\(3^a=3^b\)”的（\quad）

\choices
    {充分不必要条件}
    {必要不充分条件}
    {充要条件}
    {既不充分也不必要条件}
\end{problem}

\begin{problem}
下列图中，相关性系数最大的是（\quad）

\choices
    {\choicebitmap[width=0.15\paperwidth]{img/2024/tianjin/q3_fig1.png}}
    {\choicebitmap[width=0.15\paperwidth]{img/2024/tianjin/q3_fig2.png}}
    {\choicebitmap[width=0.15\paperwidth]{img/2024/tianjin/q3_fig3.png}}
    {\choicebitmap[width=0.15\paperwidth]{img/2024/tianjin/q3_fig4.png}}
\end{problem}

\begin{problem}
下列函数是偶函数的是（\quad）

\choices
    {\(y=\dfrac{\e^x-x^2}{\e^x+x^2}\)}
    {\(y=\dfrac{\cos x-x^2}{x^2+1}\)}
    {\(y=\dfrac{\e^x-x}{\e^x+x}\)}
    {\(y=\dfrac{\sin x-x}{x^2+1}\)}
\end{problem}

\begin{problem}
若 \(a=4.2^{-0.2}\)，\(b=4.2^{0.2}\)，\(c=\log_{4.2}0.2\)，则 \(a\)，\(b\)，\(c\) 的大小关系为（\quad）

\choices
    {\(a<b<c\)}
    {\(a<c<b\)}
    {\(c<b<a\)}
    {\(c<a<b\)}
\end{problem}

\begin{problem}
若 \(m\)，\(n\) 为两条不同的直线，\(\alpha\) 为一个平面，则下列结论中正确的是（\quad）

\choices
    {若 \(m\parallel \alpha\)，\(m\perp n\)，则 \(n\perp \alpha\)}
    {若 \(m\perp \alpha\)，\(m\perp n\)，则 \(n\perp \alpha\)}
    {若 \(m\parallel \alpha\)，\(n\perp \alpha\)，则 \(m\perp n\)}
    {若 \(m\perp \alpha\)，\(n\perp \alpha\)，则 \(m\perp n\)}
\end{problem}

\begin{problem}
已知函数 \(f(x)=3\sin \left( \omega x+\frac{\pi}{3} \right)\)（\(\omega>0\)）的最小正周期为 \(\pi\)，则 \(f(x)\) 在 \(\left[-\frac{\pi}{12},\frac{\pi}{6}\right]\) 上的最小值是（\quad）

\choices
    {\(-\dfrac{3\sqrt3}{2}\)}
    {\(-\dfrac32\)}
    {0}
    {\(\dfrac32\)}
\end{problem}

\begin{problem}
双曲线 \(\frac{x^2}{a^2}-\frac{y^2}{b^2}=1\)（\(a>0\)，\(b>0\)）的左，右焦点分别为 \(F_1\)，\(F_2\). \(P\) 是双曲线右支上一点，且直线 \(PF_2\) 的斜率为 2. \(\triangle PF_1F_2\) 是面积为 8 的直角三角形，则双曲线的方程为（\quad）

\choices
    {\(\dfrac{x^2}{2}-\dfrac{y^2}{8}=1\)}
    {\(\dfrac{x^2}{4}-\dfrac{y^2}{8}=1\)}
    {\(\dfrac{x^2}{8}-\dfrac{y^2}{2}=1\)}
    {\(\dfrac{x^2}{8}-\dfrac{y^2}{4}=1\)}
\end{problem}

\begin{problem}
一个五面体 \(ABC-DEF\). 已知 \(AD\parallel BE\parallel CF\)，且两两之间距离为 1，并已知 \(AD=1\)，\(BE=2\)，\(CF=3\). 则该五面体的体积为（\quad）

\bitmapfigure[max width=6.0cm]{img/2024/tianjin/q9_fig1.png}

\choices
    {\(\dfrac{\sqrt3}{6}\)}
    {\(\dfrac{3\sqrt3}{4}+\dfrac12\)}
    {\(\dfrac{\sqrt3}{2}\)}
    {\(\dfrac{3\sqrt3}{4}-\dfrac12\)}
\end{problem}

\section{填空题}

\begin{problem}
已知 \(\i\) 是虚数单位，复数 \((\sqrt5+\i)(\sqrt5-2\i)=\fillinblank{}\).
\end{problem}

\begin{problem}
在 \(\left(\frac{x^2}{3}+\frac{3}{x^2}\right)^6\) 的展开式中，常数项为\fillinblank{}.
\end{problem}

\begin{problem}
圆 \((x-1)^2+y^2=25\) 的圆心与抛物线 \(y^2=2px\)（\(p>0\)）的焦点 \(F\) 重合，且两曲线在第一象限的交点为 \(A\)，则原点到直线 \(AF\) 的距离为\fillinblank{}.
\end{problem}

\begin{problem}
某校组织学生参加农业实践活动，期间安排了劳动技能比赛，比赛共 5 个项目，分别为整地做畦，旱田播种，作物移栽，田间灌溉，藤架搭建，规定每人参加其中 3 个项目. 假设每人参加每个项目的可能性相同，则甲同学参加“整地做畦”项目的概率为\fillinblank{}；已知乙同学参加的 3 个项目中有“整地做畦”，则他还参加“田间灌溉”项目的概率为\fillinblank{}.
\end{problem}

\begin{problem}
在边长为 1 的正方形 \(ABCD\) 中，点 \(E\) 为线段 \(CD\) 的三等分点，\(CE=\frac12DE\)，\(\overrightarrow{BE}=\lambda\overrightarrow{BA}+\mu\overrightarrow{BC}\)，则 \(\lambda+\mu=\fillinblank{}\)；\(F\) 为线段 \(BE\) 上的动点，\(G\) 为 \(AF\) 中点，则 \(\overrightarrow{AF}\cdot \overrightarrow{DG}\) 的最小值为\fillinblank{}.

\bitmapfigure[max width=4.8cm]{img/2024/tianjin/q14_fig1.png}
\end{problem}

\begin{problem}
若函数 \(f(x)=2\sqrt{x^2-ax}-|ax-2|+1\) 恰有一个零点，则 \(a\) 的取值范围为\fillinblank{}.
\end{problem}

\section{解答题}

\begin{problem}
在 \(\triangle ABC\) 中，角 \(A\)，\(B\)，\(C\) 所对的边分别为 \(a\)，\(b\)，\(c\)，已知 \(\cos B=\frac{9}{16}\)，\(b=5\)，\(\frac{a}{c}=\frac23\).
\begin{enumerate}
    \item 求 \(a\)；
    \item 求 \(\sin A\)；
    \item 求 \(\cos(B-2A)\) 的值.
\end{enumerate}
\end{problem}

\begin{problem}
已知四棱柱 \(ABCD-A_1B_1C_1D_1\) 中，底面 \(ABCD\) 为梯形，\(AB\parallel CD\)，\(A_1A\perp\) 平面 \(ABCD\)，\(AD\perp AB\)，其中 \(AB=A_1A=2\)，\(AD=DC=1\).
\(N\) 是 \(B_1C_1\) 的中点，\(M\) 是 \(DD_1\) 的中点.
\begin{enumerate}
    \item 求证：\(D_1N\parallel\) 平面 \(CB_1M\)；
    \item 求平面 \(CB_1M\) 与平面 \(BB_1CC_1\) 的夹角余弦值；
    \item 求点 \(B\) 到平面 \(CB_1M\) 的距离.
\end{enumerate}

\bitmapfigure[max width=6.0cm]{img/2024/tianjin/q17_fig1.png}
\end{problem}

\begin{problem}
已知椭圆 \(\frac{x^2}{a^2}+\frac{y^2}{b^2}=1\)（\(a>b>0\)）的离心率 \(e=\frac12\). 左顶点为 \(A\)，下顶点为 \(B\)，\(C\) 是线段 \(OB\) 的中点，其中 \(S_{\triangle ABC}=\frac{3\sqrt3}{2}\).
\begin{enumerate}
    \item 求椭圆方程；
    \item 过点 \(\left(0,-\frac32\right)\) 的动直线与椭圆有两个交点 \(P\)，\(Q\). 在 \(y\) 轴上是否存在点 \(T\) 使得 \(\overrightarrow{TP}\cdot \overrightarrow{TQ}\le 0\)？若存在，求出这个 \(T\) 点纵坐标的取值范围；若不存在，请说明理由.
\end{enumerate}
\end{problem}

\begin{problem}
已知数列 \(\{a_n\}\) 是公比大于 0 的等比数列，其前 \(n\) 项和为 \(S_n\). 若 \(a_1=1\)，\(S_2=a_3-1\).
\begin{enumerate}
    \item 求数列 \(\{a_n\}\) 前 \(n\) 项和 \(S_n\)；
    \item 设 \(b_n=k\)（\(n=a_k\)），\(b_n=b_{n-1}+2k\)（\(a_k<n<a_{k+1}\)），其中 \(k\in \mathbb N^*\).
    \begin{enumerate}
        \item 当 \(k\ge 2\)，\(n=a_{k+1}\) 时，求证：\(b_{n-1}\ge a_k b_n\)；
        \item 求 \(\displaystyle \sum_{i=1}^{S_n}b_i\).
    \end{enumerate}
\end{enumerate}
\end{problem}

\begin{problem}
设函数 \(f(x)=x\ln x\).
\begin{enumerate}
    \item 求 \(f(x)\) 图象上点 \((1,f(1))\) 处的切线方程；
    \item 若 \(f(x)\ge a(x-\sqrt{x})\) 在 \(x\in(0,+\infty)\) 时恒成立，求 \(a\) 的值；
    \item 若 \(x_1\)，\(x_2\in(0,1)\)，证明 \(|f(x_1)-f(x_2)|\le |x_1-x_2|^{1/2}\).
\end{enumerate}
\end{problem}

